\section{Discussion}
\label{sec:discussion}

\paragraph{Limitations} \emph{A closed world.} Every result assumes that the list of candidate causes is given before the investigation starts. In real triage no one supplies this list, and the true cause is often not among those an analyst would write down in advance. When it is missing, the controller can only assign the residual \textit{other} or, worse, the listed cause the evidence fits best: in sec-triage, joint evidence names another actionable cause in 24 of 48 episodes whose cause was removed from the list (RQ5). Our finding that deciding needs no model holds in this closed world; proposing candidate causes during the investigation is plausibly where a model would contribute to deciding, and \hesp{} gives it no way to do so, since a verdict outside the list is rejected. \emph{Distinctive observations.} In sec-triage and sigma-triage almost every cause produces, with probability one, an outcome on some probe that no other cause produces, so one well-chosen probe usually settles the case. This is why a table-driven controller suffices there; in comp-triage, where causes share their observations, the single-signature rule escalates 52 of 72 episodes. The counted tables come from the same generator as the evaluation episodes, and the verifier's signatures are exactly the distinctive outcomes in those tables, so the agreement between counted and generator tables, and between the controller's stopping rule and the verifier, is expected by construction. \emph{Who wrote the settings.} All three settings are simulated, so that ground truth and the attacker can be controlled; real telemetry is noisier and larger. They are small (eight causes, twelve rules, three scenarios), and development and evaluation episodes share a generator, so our results generalize to new episodes of these tasks, not to new causes. sigma-triage reduces our influence on which causes compete, but its outcomes are a model's deterministic annotation, and one rule whose annotation lacked a distinctive outcome was replaced. With structured observations, the environment's own classifier turns each response into an outcome with fixed rules; the model's reading is tested only on raw logs in two formats, the drifted one written by us. \emph{Given and exclusive causes.} The causes are also mutually exclusive, and the environment reports when an incident changes; concurrent causes and unreported change are outside our scope. \emph{Our attackers.} We tested four injected texts, two forgeries, and, against the reader, one changed format, one claim wording, and one log-line forgery. An attacker who controls two independent source groups defeats corroboration, and reader trust is only as strong as the rule parser it trusts; ours is immune to the tested attacks by construction, since it reads the first record of each probe. \emph{Two small models.} The prompt comparison and the reader attacks ran on Qwen2.5-7B and Llama-3.1-8B only. \emph{Strict evidence.} Our verifier rejects correct verdicts reached by combining observations; \S\ref{sec:policy} measures what this costs, but the comp-triage scenarios that do so are few and written by us.

\paragraph{Attacks beyond our threat model} Three attacks lie outside the goals \hesp{} meets. An attacker who controls two independent source groups can satisfy corroboration with forged evidence; raising the number of required groups raises the bar but escalates more honest benign cases, as RQ2 shows for causes with a single signature. An attacker who poisons the development cases from which the tables are counted shifts every future verdict; tables should be counted from cases whose resolution does not depend on the agent. Finally, an attacker can suppress a signature instead of forging one, for example by avoiding the behavior a probe detects; this turns an attack into an escalation rather than a benign closure, because \hesp{} never accepts a verdict by elimination.

\paragraph{Deployment} We recommend \hesp{} with information-gain selection, the finish guard, the confirmation stop, source-diverse corroboration, and recovery, and, wherever the model parses raw logs, a rule parser first with reader trust on top. Verdicts should reach an analyst with their cited evidence and journal, so that a residual attack becomes a delayed response rather than an unattended closure. Two inputs determine how well this works. The source groups must be declared from the deployment's actual data flow, since corroboration is only as independent as they are. The prediction tables must come from cases that resemble production; resolved tickets record only the probes an analyst chose, so tables estimated from them face selection bias, and the confirmation stop additionally needs, for each cause, at least one probe that singles it out.

\paragraph{The role of the model} With structured observations and a given list of causes, the controller accounts for all of the measured gain, and the advantage that strong models once showed came from waiting for a confirming signature, which a rule reproduces. The model is needed to read logs whose format a fixed parser cannot handle, and that is exactly where it is dangerous. The component to harden in a local triage agent is therefore the reader, not the planner. Two further roles lie outside our settings and are where we would look next: proposing candidate causes that the list lacks, and weighing evidence that only adds up in combination, as in comp-triage. Both require settings that a fixed procedure cannot solve, and the controller alone is the baseline any such study must beat.

\paragraph{Public security data}\label{sec:realdata} To test these results on real data, we asked whether the three public datasets that come closest could evaluate an investigation policy, which needs alerts that leave several explanations plausible, probes that differ in cost and informativeness, and objective ground truth (Appendix~\ref{app:public}). In the subsets we examined, none could: most ExCyTIn-Bench~\citep{excytin} questions name the alert to locate; in GUIDE~\citep{guide}, organization identity carries 1.03 of the grade's 1.50 bits; and in the OTRF~\citep{otrf} remote-execution recordings we audited, a rule on the alert text separates 44 of 46 alerts. What these subsets lack are paired attacks and legitimate administration that use the same tools, under objective labels.
